\subsection{Marker-ring core}

\begin{theorem}[Jeandel--Rao marker-ring $b_3$ core]\label{thm:jeandel-rao-b3}\lean{Obsmetrologia.AperiodicOrder.Tiling.Witness.markerRing_b3_nonzero_h1_class}\leanok\uses{thm:jeandel-rao-witness}
The logically prior certificate the decorated witness is pulled back from: in
the marker-ring quotient of the $444$ collared Jeandel--Rao states, the quotient
has $4$ faces $g_0,\ldots,g_3$ and $7$ edges $q_0,\ldots,q_6$. The three
nonzero boundary columns are dependent,
$\partial g_0+\partial g_2+\partial g_3=0$, so
$\operatorname{rank}\partial=2$ and $H^1$ has rank $5$. The class
$b_3=q_1+q_2+q_5$ is an isolated nonzero cocycle. This is still a finite
witness only, but it is the marker-ring core from which the decorated cochain
is pulled back.
\end{theorem}
